\documentclass[14pt,a4paper]{extreport}
% preamble.tex — спільна преамбула для всіх лабораторних робіт. Підключається одразу після \documentclass у кожному lab*/main.tex.
%
% ЗБІРКА: LuaLaTeX (не pdfLaTeX). Причина: пакет listings під pdfLaTeX не читає UTF-8 і зупиняє збірку на кириличних коментарях у лістингах програмного коду. LuaLaTeX обробляє UTF-8 нативно.
%
% ШРИФТ: Liberation Serif — метрично сумісний з Times New Roman і містить повний набір кириличних гліфів (TeX Gyre Termes кирилиці не має, через що український текст не відображався взагалі).

\usepackage{fontspec}
\setmainfont{Liberation Serif}
\setmonofont{DejaVu Sans Mono}[Scale=0.82]

\usepackage{polyglossia}
\setmainlanguage{ukrainian}
\setotherlanguage{english}

\usepackage{microtype}

% Поля, інтервал, відступи й заголовки нижче взяті з шаблону НУЛП, за яким написана бакалаврська робота (../bachelor-thesis/thesis/main.tex). Той самий формат приймає В.М. Шиманський і для магістерських, тож усі звіти тримаємо в ньому, а не в тому, що переказували на консультації.
\usepackage[left=25mm,right=15mm,top=15mm,bottom=23mm,footskip=1.5mm]{geometry}
\usepackage{setspace}
% Саме \setstretch{1.5}, а не \onehalfspacing: останній дає близько 1.25.
\setstretch{1.5}

\usepackage{graphicx}
\usepackage{float}
\usepackage{booktabs}
\usepackage{array}
\usepackage{tabularx}
\usepackage{longtable}
\usepackage{enumitem}
% labelindent вирівнює маркер списку з абзацним відступом — так у шаблоні НУЛП.
\setlist{nosep}
\setlist[itemize]{label=\textbullet,leftmargin=*,labelindent=1.27cm}
\setlist[enumerate]{leftmargin=*,labelindent=1.27cm}
\usepackage{indentfirst}
\usepackage{amsmath}

\usepackage{hyperref}
\hypersetup{colorlinks=true,linkcolor=black,urlcolor=blue,citecolor=black}

% xurl дозволяє \url розривати адресу в будь-якому місці, а не лише після скісної риски. Без нього довгий сегмент шляху (напр. «scopus---podrobnaya- instruktsiya-dlya-uchenogo») виходить за межі сторінки, бо дефіс для hyperref не є точкою розриву. Вантажиться після hyperref — перевизначає його \UrlBreaks.
\usepackage{xurl}

% ---------------------------- Лістинги програмного коду ---------------------------- Використано minted (через Pygments), а не listings: listings не вміє коректно обробляти багатобайтові символи й «висмикує» кириличні коментарі з коду, зліплюючи їх в один рядок. Потребує прапорця -shell-escape.
\usepackage{xcolor}
\usepackage[newfloat]{minted}
% labelsep=period дає «Рис. 2.1. Назва» замість «Рис. 2: Назва» — формат НУЛП.
\usepackage[labelsep=period]{caption}

% baselinestretch=1 обов'язковий: без нього лістинг успадковує півторашку основного тексту, і довгий фрагмент коду перестає вміщатися у сторінку. Усередині [H]-флоата це не помилка збірки, а нескінченний цикл порожніх сторінок — lab1 так намолотив 161 тисячу сторінок.
\setminted{
  baselinestretch=1,
  fontsize=\footnotesize,
  linenos=true,
  frame=single,
  framesep=4pt,
  breaklines=true,
  breakanywhere=true,
  tabsize=4,
  numbersep=7pt,
  bgcolor=gray!8,
}
\setminted[python]{python3=true}

\SetupFloatingEnvironment{listing}{name=Лістинг,listname=Перелік лістингів}

% ---------------------------- Нумерація та оформлення заголовків ----------------------------
\usepackage{titlesec}

\setcounter{secnumdepth}{3}
\renewcommand{\thesection}{\arabic{section}}
\renewcommand{\thesubsection}{\arabic{section}.\arabic{subsection}}
\renewcommand{\thesubsubsection}{\arabic{section}.\arabic{subsection}.\arabic{subsubsection}}

% Розділ по центру, 16 пт, великими літерами; слово «РОЗДІЛ» опущене, бо звіт не дисертація. Номер лишається: \numberwithin прив'язує до нього нумерацію рисунків («Рис. 2.1»), і без нумерованих розділів вона б зламалася. \MakeUppercase четвертим аргументом отримує текст заголовка, тому в джерелі заголовки пишуться звичайним регістром.
\titleformat{\section}[block]
  {\centering\bfseries\fontsize{16}{19.2}\selectfont}
  {\thesection.}{0.3em}{\MakeUppercase}
% Підрозділи — зліва з абзацного відступу, звичайним регістром.
\titleformat{\subsection}[block]{\bfseries}{\thesubsection.}{1em}{}
\titleformat{\subsubsection}[block]{\bfseries\itshape}{\thesubsubsection.}{1em}{}

% Незіркові \titlespacing — обов'язкові. Без них titlesec застосовує зірковий варіант, який прибирає абзацний відступ у першого абзацу після заголовка й перекриває indentfirst. Значення — типові для report.
\titlespacing{\section}{0pt}{20pt}{6pt}
\titlespacing{\subsection}{1.27cm}{6pt}{0pt}
\titlespacing{\subsubsection}{1.27cm}{6pt}{0pt}

\setlength{\parindent}{1.27cm}
\setlength{\parskip}{0pt}
\renewcommand{\arraystretch}{1.3}

% Переноси слів вимкнені — так у шаблоні НУЛП. Натомість абзац розганяється міжслівними пробілами, тому tolerance і emergencystretch підняті: без них рядки виходили б за праве поле.
\hyphenpenalty=10000
\exhyphenpenalty=10000
\tolerance=9999
\setlength{\emergencystretch}{3em}

% ---------------------------- Підписи до рисунків і таблиць ----------------------------
\captionsetup[figure]{name=Рис.,justification=centering,font={bf,it}}
\captionsetup[table]{name=Таблиця,justification=raggedleft,singlelinecheck=false,labelfont=bf,textfont=normalfont}

% Рисунки, таблиці й формули нумеруються в межах розділу: «Рис. 2.1», а не «Рис. 5». Так у шаблоні НУЛП, і так само виглядали приклади на консультації.
\numberwithin{equation}{section}
\numberwithin{figure}{section}
\numberwithin{table}{section}

% ---------------------------- Колонтитули ----------------------------
\usepackage{fancyhdr}
\pagestyle{fancy}
\fancyhf{}
\fancyfoot[R]{\thepage}
\renewcommand{\headrulewidth}{0pt}
\fancypagestyle{plain}{%
  \fancyhf{}
  \fancyfoot[R]{\thepage}
  \renewcommand{\headrulewidth}{0pt}
}

% titlepage.tex — титульний аркуш, спільний для всіх дисциплін.
% Потребує визначених у subject.tex команд:
%   \SubjectName, \LecturerName, \StudentName, \StudentGroup, \AcademicYear
%
% Використання у тексті: \MakeLabTitle{1}{Тема роботи}

\newcommand{\MakeLabTitle}[2]{%
\begin{titlepage}
\begin{center}
  {\large\bfseries МІНІСТЕРСТВО ОСВІТИ І НАУКИ УКРАЇНИ}\\[6pt]
  {\large\bfseries Національний університет «Львівська політехніка»}\\[6pt]
  {\large\bfseries Інститут комп'ютерних наук та інформаційних технологій}\\[6pt]
  {\large\bfseries Кафедра «Системи штучного інтелекту»}\\[1.2cm]

  \includegraphics[width=5cm]{../../../../logo}\\[1.2cm]

  {\Large\bfseries ЛАБОРАТОРНА РОБОТА №#1}\\[0.5cm]
  з дисципліни\\[4pt]
  «\SubjectName»\\[0.6cm]
  на тему\\[4pt]
  \textit{#2}\\[2cm]
\end{center}

\begin{flushright}
  {\bfseries Виконав:}\\
  студент групи \StudentGroup\\
  \StudentName\\[1cm]
  {\bfseries Викладач:}\\
  \LecturerName\\
\end{flushright}

\vfill
\begin{center}
  Львів --- \AcademicYear
\end{center}
\end{titlepage}
\clearpage
}


% Метадані дисципліни. Використовуються титульним аркушем (titlepage.tex).

\newcommand{\SubjectName}{Штучний інтелект та інноваційні інформаційні технології}
% Два рядки: посада з науковим ступенем, далі прізвище. `\\` розриває рядок усередині flushright-блоку титульного аркуша.
\newcommand{\LecturerName}{доц. каф. СШІ к.т.н.,\\ Думин І.Б.}
\newcommand{\StudentName}{Галич Максим}
\newcommand{\StudentGroup}{КНСШ-13}
\newcommand{\AcademicYear}{2026}

% Рисунки — згенеровані артефакти, тому лежать не поряд із текстом,
% а в .cache/ (див. ARTIFACTS у src/config.py). Імена всередині
% (eda/..., lab1/...) не змінилися.
\graphicspath{{../../../../.cache/reports-artifacts/1th term 2026 Autumn/Штучний інтелект та інноваційні інформаційні технології/}}

\begin{document}

\MakeLabTitle{1}{Побудова та аналіз моделей штучного інтелекту з використанням TensorFlow}

\textbf{Мета роботи.} Опанувати практику створення й оптимізації моделей
машинного навчання у TensorFlow, навчитися застосовувати різні архітектури,
регуляризацію та методи покращення узагальнюючої здатності моделі, виконати
порівняльний аналіз результатів.

\section{Вибір набору даних}

Для роботи обрано набір даних \textbf{California Housing} (перепис населення
США 1990 року), доступний через \texttt{sklearn.datasets.fetch\_california\_housing}.
Набір містить \textbf{20\,640} записів і \textbf{8} числових ознак; цільова
змінна \texttt{MedHouseVal} --- медіанна вартість будинку в блоці, виміряна
в сотнях тисяч доларів. Таким чином, розв'язується \textbf{задача регресії}.

\begin{table}[H]
\caption{Ознаки набору даних California Housing}
\label{tab:features}
\begin{tabularx}{\linewidth}{|l|X|}
\hline
\textbf{Ознака} & \textbf{Зміст} \\
\hline
MedInc & Медіанний дохід домогосподарств у блоці \\
\hline
HouseAge & Медіанний вік будинків \\
\hline
AveRooms & Середня кількість кімнат на домогосподарство \\
\hline
AveBedrms & Середня кількість спалень на домогосподарство \\
\hline
Population & Населення блоку \\
\hline
AveOccup & Середня кількість мешканців на домогосподарство \\
\hline
Latitude & Географічна широта \\
\hline
Longitude & Географічна довгота \\
\hline
\end{tabularx}
\end{table}

\section{Розвідувальний аналіз даних}

\subsection{Розподіл ознак}

Гістограми розподілу всіх ознак та цільової змінної наведено на
рисунку~\ref{fig:dist}. Ознаки \texttt{AveRooms}, \texttt{AveBedrms},
\texttt{Population} та \texttt{AveOccup} мають виражені довгі «хвости»
з поодинокими екстремальними значеннями, тоді як \texttt{Latitude} і
\texttt{Longitude} демонструють двомодальний розподіл, що відповідає двом
великим агломераціям штату --- Лос-Анджелесу та району затоки Сан-Франциско.

\begin{figure}[H]
  \centering
  \includegraphics[width=\linewidth]{eda/distributions}
  \caption{Розподіл ознак та цільової змінної}
  \label{fig:dist}
\end{figure}

\subsection{Кореляційний аналіз}

Кореляційну матрицю Пірсона зображено на рисунку~\ref{fig:corr}. Найсильніший
зв'язок із цільовою змінною має медіанний дохід: $r(\texttt{MedInc},
\texttt{MedHouseVal}) = 0{,}689$. Решта ознак корелює з ціною значно слабше
(за модулем не більше 0{,}152). Окремо варто відзначити дуже високу кореляцію
між \texttt{AveRooms} та \texttt{AveBedrms} ($r = 0{,}848$), що вказує на
мультиколінеарність і додатково обґрунтовує застосування методу головних
компонент.

\begin{figure}[H]
  \centering
  \includegraphics[width=0.85\linewidth]{eda/correlation}
  \caption{Кореляційна матриця ознак та цільової змінної}
  \label{fig:corr}
\end{figure}

\subsection{Аналіз дисбалансу}

Оскільки розв'язується задача регресії, класового дисбалансу в прямому
розумінні не існує. Натомість проаналізовано два споріднені ефекти, які
так само здатні зміщувати навчання (рисунок~\ref{fig:imbalance}).

По-перше, розподіл цільової змінної є \textbf{правосторонньо асиметричним}:
коефіцієнт асиметрії складає $0{,}978$. По-друге, виявлено \textbf{штучне
обмеження зверху}: рівно 992 записи (4{,}81~\% вибірки) мають значення
цільової змінної 5{,}00001, тобто їхню справжню вартість було «обрізано»
на межі 500 тис. доларів ще під час формування набору даних. Ці записи
принципово не можуть бути передбачені точно жодною моделлю і формують
нижню межу досяжної похибки.

Для коректного врахування перекосу цільову змінну розбито на 5 квантильних
груп, за якими виконано \textbf{стратифікований} поділ вибірки. Завдяки цьому
частка дорогих та дешевих об'єктів у тренувальній, валідаційній і тестовій
вибірках збігається.

\begin{figure}[H]
  \centering
  \includegraphics[width=\linewidth]{eda/target_imbalance}
  \caption{Перекіс розподілу цільової змінної та ефект обмеження зверху}
  \label{fig:imbalance}
\end{figure}

\section{Попередня обробка даних}

\subsection{Стандартизація}

Ознаки мають принципово різні масштаби (\texttt{Population} сягає десятків
тисяч, тоді як \texttt{AveBedrms} коливається близько одиниці), тому
застосовано \texttt{StandardScaler}. Критично важливо, що параметри
масштабування оцінюються \textbf{виключно на тренувальній вибірці} і лише
потім застосовуються до валідаційної та тестової --- інакше статистики
тестових даних просочилися б у процес навчання (data leakage).

\subsection{Метод головних компонент}

Результат PCA наведено на рисунку~\ref{fig:pca}. Для збереження 90~\%
дисперсії потрібно \textbf{5 з 8} компонент. Перша компонента пояснює лише
25{,}3~\% дисперсії, тобто дані не мають однієї домінуючої осі мінливості,
а інформація розподілена між ознаками досить рівномірно. З огляду на це
для навчання моделей використано \textbf{повний набір з 8 ознак}: відкидання
трьох компонент коштувало б 10~\% дисперсії, не даючи помітного виграшу
у швидкості на вибірці такого розміру.

Додатково виконано нелінійну проєкцію t-SNE на випадковій підвибірці
з 3000 записів. Проєкція показує, що об'єкти з близькою вартістю формують
зв'язні області, тобто ознаковий простір дійсно несе інформацію, достатню
для регресії.

\begin{figure}[H]
  \centering
  \includegraphics[width=\linewidth]{eda/pca_tsne}
  \caption{PCA: накопичена пояснена дисперсія та t-SNE-проєкція}
  \label{fig:pca}
\end{figure}

\subsection{Поділ вибірки}

Вибірку поділено у співвідношенні \textbf{60/20/20} зі стратифікацією за
квантильними групами цільової змінної: 12\,384 записи для навчання,
4\,128 для валідації та 4\,128 для тестування.

\begin{listing}[H]
\inputminted[firstline=32,lastline=56]{python}{../src/aiit/data.py}
\caption{Поділ вибірки та масштабування (\texttt{src/data.py})}
\label{lst:data}
\end{listing}

\section{Архітектури нейронних мереж}

Відповідно до завдання реалізовано дві архітектури, кожну --- у двох
варіантах функції активації (ReLU та GELU), що дає чотири конфігурації.

\begin{table}[H]
\caption{Реалізовані конфігурації моделей}
\label{tab:arch}
\begin{tabularx}{\linewidth}{|l|X|r|}
\hline
\textbf{Модель} & \textbf{Склад} & \textbf{Параметрів} \\
\hline
baseline & Dense(64) → Dense(32) → Dense(1), без регуляризації & 2\,689 \\
\hline
regularized & [Dense → BatchNorm → Activation → Dropout(0,2)] ×2 → Dense(1),
L2 $\lambda = 10^{-4}$ & 2\,977 \\
\hline
\end{tabularx}
\end{table}

\begin{listing}[H]
\inputminted[firstline=20,lastline=50]{python}{../src/aiit/tf_models.py}
\caption{Архітектури моделей у TensorFlow (\texttt{src/tf\_models.py})}
\label{lst:tfmodels}
\end{listing}

\section{Навчання моделей}

Усі моделі навчалися з однаковими гіперпараметрами: оптимізатор \textbf{Adam},
швидкість навчання $10^{-3}$, розмір батча 256, \textbf{20 епох}, функція
втрат MSE. Фіксовані зерна генераторів випадкових чисел забезпечують
відтворюваність.

\begin{listing}[H]
\inputminted[firstline=51,lastline=89]{python}{../src/aiit/run_lab1.py}
\caption{Цикл навчання та оцінювання (\texttt{src/run\_lab1.py})}
\label{lst:run1}
\end{listing}

\section{Результати}

\subsection{Метрики якості}

Оскільки задача регресійна, використано метрики RMSE, MAE, коефіцієнт
детермінації $R^2$ та MAPE. RMSE і MAE виражені в тих самих одиницях, що
й цільова змінна (сотні тисяч доларів), тому допускають пряму інтерпретацію.

\begin{table}[H]
\caption{Метрики моделей на тестовій вибірці (TensorFlow)}
\label{tab:tfresults}
\begin{tabularx}{\linewidth}{|X|r|r|r|r|r|}
\hline
\textbf{Модель} & \textbf{RMSE} & \textbf{MAE} & \textbf{$R^2$} &
\textbf{MAPE, \%} & \textbf{Час, с} \\
\hline
baseline-relu       & \textbf{0,6082} & \textbf{0,4257} & \textbf{0,7191} & 23,7 & 7,60 \\
\hline
baseline-gelu       & 0,6181 & 0,4301 & 0,7098 & 24,2 & 2,29 \\
\hline
regularized-relu    & 0,7480 & 0,4928 & 0,5751 & 23,3 & 6,84 \\
\hline
regularized-gelu    & 0,7770 & 0,5126 & 0,5415 & 23,6 & 8,04 \\
\hline
\end{tabularx}
\end{table}

Найкращий результат показала базова мережа з активацією ReLU:
RMSE $= 0{,}608$, що відповідає середній похибці близько 60 тис. доларів,
та $R^2 = 0{,}719$ --- модель пояснює майже 72~\% дисперсії вартості житла.

\subsection{Навчальні криві}

Навчальні криві для всіх чотирьох конфігурацій наведено на
рисунку~\ref{fig:tfcurves}.

\begin{figure}[H]
  \centering
  \includegraphics[width=\linewidth]{lab1/learning_curves}
  \caption{TensorFlow: криві тренувальної та валідаційної втрати}
  \label{fig:tfcurves}
\end{figure}

\subsection{Діагностика перенавчання}

Різницю між валідаційною та тренувальною втратою зображено на
рисунку~\ref{fig:tfgap}. На другій епосі регуляризовані моделі дають
великий сплеск розриву (0,79 та 0,90), однак він швидко спадає: це
артефакт шару BatchNorm, ковзні статистики якого на початку навчання ще
погано оцінені, через що валідаційна втрата тимчасово завищена.

Починаючи приблизно з сьомої епохи розрив для всіх чотирьох моделей
стабілізується і наприкінці навчання не перевищує \textbf{0,035}, причому
не зростає з часом. Отже, \textbf{перенавчання відсутнє}.

\begin{figure}[H]
  \centering
  \includegraphics[width=0.8\linewidth]{lab1/overfit_gap}
  \caption{TensorFlow: розрив між валідаційною та тренувальною втратою}
  \label{fig:tfgap}
\end{figure}

Натомість спостерігається протилежний ефект --- \textbf{недонавчання
регуляризованих моделей}. Їхня втрата висока \textit{і на тренувальній,
і на валідаційній вибірці} (0,520 проти 0,336 у базової моделі з ReLU), що
є класичною ознакою надмірної регуляризації, а не перенавчання.

\subsection{Якість прогнозу найкращої моделі}

\begin{figure}[H]
  \centering
  \includegraphics[width=0.62\linewidth]{lab1/pred_vs_actual}
  \caption{TensorFlow: прогнозовані значення проти фактичних (baseline-relu)}
  \label{fig:tfpred}
\end{figure}

На рисунку~\ref{fig:tfpred} добре помітна вертикальна «смуга» точок при
фактичному значенні 5,0 --- це саме ті 992 записи з обмеженою зверху
вартістю, для яких модель прогнозує широкий діапазон значень. Вони дають
істотний внесок у підсумкову похибку, і усунути його засобами моделювання
неможливо.

\section{Висновки}

У ході лабораторної роботи засобами TensorFlow побудовано та проаналізовано
чотири конфігурації нейромережевих моделей для задачі регресії на наборі
даних California Housing.

\begin{enumerate}
\item Розвідувальний аналіз показав, що визначальною ознакою є медіанний
дохід ($r = 0{,}689$), а розподіл цільової змінної асиметричний
($0{,}978$) та штучно обмежений зверху для 4,81~\% записів. Останнє формує
нижню межу досяжної похибки.

\item Застосовано стандартизацію ознак і стратифікований за квантилями
поділ вибірки 60/20/20. Аналіз PCA показав, що 90~\% дисперсії потребує
5 з 8 компонент, тому зменшення розмірності визнано недоцільним.

\item Найкращою виявилася \textbf{базова} мережа з активацією ReLU:
RMSE $= 0{,}608$, $R^2 = 0{,}719$.

\item Регуляризація \textbf{погіршила} результат: RMSE зросла з 0,608 до
0,748 ($+23$~\%). Причина --- недонавчання: комбінація Dropout 0,2,
BatchNorm та L2 надто агресивна для компактної мережі з 2,7 тис. параметрів
на 12 тис. тренувальних прикладів. Регуляризація виправдана там, де наявне
перенавчання; у цьому експерименті його не було, тому вона лише обмежила
здатність моделі апроксимувати залежність.

\item Порівняння функцій активації показало незначну перевагу ReLU над GELU
для базової архітектури (RMSE 0,608 проти 0,618). Для цієї задачі вибір
активації впливає на результат значно слабше, ніж вибір архітектури.
\end{enumerate}

\end{document}
